\documentclass[UTF8,11pt]{ctexart}
\usepackage[a4paper,margin=22mm]{geometry}
\usepackage{amsmath,amssymb,bm,graphicx,adjustbox,fvextra,longtable,array}
\xeCJKDeclareCharClass{CJK}{"2160 -> "217F, "2460 -> "249B, "2200 -> "22FF}
\newcommand{\fitmath}[1]{\begin{adjustbox}{max width=\linewidth}$\displaystyle #1$\end{adjustbox}}
\setlength{\parindent}{0pt}\setlength{\parskip}{.65em}\renewcommand{\arraystretch}{1.3}
\sloppy\allowdisplaybreaks
\begin{document}
\section*{2023年考研政治答案解析}
2023年全国硕士研究生


\subsection*{招生考试思想政治理论试题详解}

\subsection*{一、单项选择题}

\subsection*{1. 【答案】A}

\subsection*{【考点知识还原】}

见表8


\subsection*{【考点点拨】}

历史思维能力


\subsection*{【试题简析】}

历史思维能力是辩证思维与历史眼光的结合，是马克思主义科学历史观的具体表现和实践运用，是以史为鉴、知古鉴今，善于运用历史眼光认识发展规律、把握前进方向、指导现实工作的能力。历史、现实、未来是相通的，历史是过去的现实，现实是未来的历史。历史思维能力的培养，能够使人正确理解和掌握历史知识，认识历史发展规律，进而对社会现实问题进行科学的观察与思考。培养并不断提高历史思维能力，是马克思主义科学世界观和方法论的内在要求。“勿忘昨天的苦难辉煌，无愧今天的使命担当，不负明天的伟大梦想。”党的十九届六中全会公报中的这句话意味深长，既接续历史又映照现实、指引未来。回顾总结百年奋斗光辉历程，昨天、今天、明天紧密相连，充分体现了唯物史观、大历史观和正确党史观。“编纂出版《复兴文库》大型历史文献丛书，就是要通过对近代以来重要思想文献的选编，述录先人的开拓，启迪来者的奋斗。”“历史是最好的教科书，一切向前走，都不能忘记走过的路；走得再远、走到再光辉的未来，也不能忘记走过的过去。”这表明历史、现实、未来是相通的，A 正确，B不符合题意且内容错误。历史过程既包含必然的因素，也包含偶然的因素。历史不仅存在着起决定作用的一般规律或者历史必然性，而且还存在着影响历史进程的诸多偶然性。马克思指出：“如果斗争只是在机会绝对有利的条件下才着手进行，那么创造世界历史未免就太容易了。另一方面，如果‘偶然性’不起任何作用的话，那么世界历史就会带有非常神秘的性质。这些偶然性本身自然纳入总的发展过程中，并且为其他偶然性所补偿。但是，发展的加速和延缓在很大程度上是取决于这些‘偶然性’的，其中也包括一开始就站在运动最前面的那些人物的性格这样一种‘偶然情况’。”C错误。社会历史是有规律的，但是，历史事件不可能完整重复和再现，因为每一个具体的历史事件都是历史主、客体多种复杂因素相互作用的结果，是在一定时间、地点、人物等条件下发生的，它总是个别的、具体的、不可重复的，D错误。


\subsection*{2. 【答案】C}

\subsection*{【考点知识还原】}

见表 17


\subsection*{【考点点拨】}

划分社会形态的依据


\subsection*{【试题简析】}

社会形态是关于社会运动的具体形式、发展阶段和不同质态的范畴，是同生产力发展一定阶段相适应的经济基础与上层建筑的统一体。经济基础是指由社会一定发展阶段的生产力所决定的生产关系的总和。生产关系是人们在物质生产过程中形成的不以人的意志为转移的经济关系，是社会关系中最基本的关系。政治关系、家庭关系、宗教关系等其他社会关系，都受生产关系的支配和制约。生产关系包括生产资料所有制关系、生产中人与人的关系和产品分配关系。在生产关系中，生产资料所有制关系是最基本的、决定性的，它是人们进行物质资料生产的前提，生产、分配、交换和消费关系在很大程度上是由这一前提决定的。它是区分不同生产方式、判定社会经济结构性质的客观依据。C正确。A、D属于生产力范畴，B属于上层建筑范畴，都不是判定、划分社会形态的依据。


\subsection*{3. 【答案】C}

\subsection*{【考点知识还原】}

见表21


\subsection*{【考点点拨】}

商品价值量的决定


\subsection*{【试题简析】}

商品的价值包括质的规定与量的规定两个方面。价值的质的规定回答的是价值的实体是什么，价值的量的规定则回答价值的大小由什么决定和怎样决定。商品的价值是凝结在商品中的无差别的一般人类劳动，价值量是由劳动者生产商品所耗费的劳动量决定的，而劳动量则按照劳动时间来计量。决定商品价值量的不是生产商品的个别劳动时间，而是社会必要劳动时间。生产商品所需要的社会必要劳动时间随着劳动生产率的变化而变化。劳动生产率指的是劳动者生产使用价值的效率。它的高低可以用单位劳动时间内生产的产品数量来测量，也可以用单位产品中所耗费的劳动时间来测量。劳动生产率水平越高，单位时间内生产的商品数量越多，生产每件商品所需要的社会必要劳动时间就越少，单位商品的价值量就越小；反之就越大。商品的价值量与生产商品所耗费的社会必要劳动时间成正比，与劳动生产率成反比。C正确，A、B错误。不管劳动生产率怎么变化，同一劳动在同样的时间内提供的价值总量是不变的，D错误。


\subsection*{4. 【答案】A}

\subsection*{【考点知识还原】}

见表30


\subsection*{【考点点拨】}

垄断组织的获利


\subsection*{【试题简析】}

垄断利润主要是通过垄断组织制定的垄断价格来实现的。垄断价格是垄断组织在销售或购买商品时，凭借其垄断地位规定的、旨在保证获取最大限度利润的市场价格。垄断价格包括垄断高价和垄断低价两种形式。垄断高价是指垄断组织出售商品时规定的高于生产价格的价格；垄断低价是指垄断组织在购买非垄断企业所生产的原材料等生产资料时规定的低于生产价格的价格。本题首先可以排除C、D（垄断低价和垄断高价）。在垄断资本主义条件下，垄断组织操纵价格带来的结果是抑制了市场上价格的自由波动，垄断价格一定时期内背离生产价格和价值，价值规律是通过垄断价格的变动来实现的。由于垄断地位产生的垄断价格，不能长期地过分高于产品的价值或生产价格，它也会受到产品成本和商品供求的制约。垄断利润是垄断价格与生产成本之间的差额。由于市场波动，垄断价格也是会低于生产价格的，为了避免遭受净亏损，垄断组织所能接受的最低价格是成本价格（生产成本是由生产中实际耗费的不变资本和可变资本所构成）。A正确。生产价格是成本价格加平均利润，若以生产价格出售，垄断组织还是可以获得平均利润的，这与题干所设定的为了避免遭受“净亏损”不符，B错误。


\subsection*{5. 【答案】D}

\subsection*{【考点知识还原】}

见表44


\subsection*{【考点点拨】}

搞清楚什么是社会主义、怎样建设社会主义的关键


\subsection*{【试题简析】}

什么是社会主义、怎样建设社会主义，是邓小平在领导改革开放和现代化建设这一新的革命过程中，不断提出和反复思考的首要的基本的理论问题。我国社会主义在改革开放前所经历的曲折和失误，归根到底就在于对这个问题没有完全搞清楚；改革开放以来在前进中遇到的一些犹疑和困惑，归根到底也在于对这个问题没有完全搞清楚。搞清楚什么是社会主义、怎样建设社会主义，关键是要在坚持社会主义基本制度的基础上进一步认清社会主义的本质。邓小平根据马克思主义的基本原理和社会主义的实践经验，对这个问题进行了不懈的探索，作出了科学的回答。D正确。


\subsection*{【名师点拨】}

本题与 2020 年第5题考查的是同一个知识点。


\subsection*{6. 【答案】C}

\subsection*{【考点知识还原】}

见表63


\subsection*{【考点点拨】}

构建初次分配、再分配、第三次分配协调配套的制度体系


\subsection*{【试题简析】}

习近平在《正确认识和把握我国发展重大理论和实践问题》中指出，要发挥分配的功能和作用，处理好效率和公平关系，构建初次分配、再分配、第三次分配协调配套的基础性制度安排。要坚持按劳分配为主体，提高劳动报酬在初次分配中的比重，完善按要素分配政策。要发挥再分配的调节作用，加大税收、社保、转移支付等的调节力度，提高精准性。要发挥好第三次分配作用，引导、支持有意愿有能力的企业和社会群体积极参与公益慈善事业，但不能搞道德绑架式“逼捐”。这里，初次分配即第一次收入分配。初次分配注重效率，是按社会生产要素（如劳动、资本、土地等）的贡献进行的分配；再分配是在初次分配结果的基础上，政府通过税收、社保、转移支付等，使各收入主体之间实现现金或实物转移的一种收入再次分配过程，也是对要素收入进行再次调节的过程；第三次分配是建立在自愿性的前提下，以募资、自行捐助等公益慈善方式对社会资源和物质财富进行的分配。C正确，A、B、D错误。


\subsection*{7. 【答案】D}

\subsection*{【考点知识还原】}

见表 70


\subsection*{【考点点拨】}

总体国家安全观


\subsection*{【试题简析】}

实现中华民族伟大复兴是中华民族近代以来最伟大的梦想。现在，我们比历史上任何时期都更接近中华民族伟大复兴的目标，比历史上任何时期都更有信心、有能力实现这个目标。同时，必须清醒地看到，我们越发展壮大，遇到的阻力和压力就会越大。我们要开创中华民族伟大复兴新局面，就必须冷静审视复杂变化的国际形势，牢固树立总体国家安全观，坚持国家利益至上，以人民安全为宗旨，以政治安全为根本，以经济安全为基础，以军事、科技、文化、社会安全为保障，以促进国际安全为依托，统筹外部安全和内部安全、国土安全和国民安全、传统安全和非传统安全、自身安全和共同安全，统筹维护和塑造国家安全，时刻准备应对重大挑战、抵御重大风险，不断把中华民族伟大复兴事业推向前进。坚持总体国家安全观，归根到底是为实现中华民族伟大复兴提供坚强安全保障，确保中华民族伟大复兴进程不被迟滞甚至中断。D正确。


\subsection*{8. 【答案】B}

\subsection*{【考点知识还原】}

见表69


\subsection*{【考点点拨】}

自我革命


\subsection*{【试题简析】}

党的十八大以来，习近平多次谈到如何跳出治乱兴衰的历史周期率的问题。2021年11月11日，习近平在党的十九届六中全会第二次全体会议上的讲话中指出，我们党历史这么长、规模这么大、执政这么久，如何跳出治乱兴衰的历史周期率？毛泽东同志在延安窑洞里给出了第一个答案，这就是“只有让人民来监督政府，政府才不敢松懈”。经过百年奋斗特别是党的十八大以来新的实践，我们党又给出了第二个答案，这就是自我革命。在党的二十大报告中，习近平强调：“经过不懈努力，党找到了自我革命这一跳出治乱兴衰历史周期率的第二个答案，自我净化、自我完善、自我革新、自我提高能力显著增强，管党治党宽松软状况得到根本扭转，风清气正的党内政治生态不断形成和发展，确保党永远不变质、不变色、不变味。”B正确。


\subsection*{9. 【答案】A}

\subsection*{【考点知识还原】}

见表73


\subsection*{【考点点拨】}

导致中国落后于时代的原因


\subsection*{【试题简析】}

中华民族是世界上古老而伟大的民族，创造了绵延5000 多年加烂文明，为人类文明进步作出了不可磨灭的贡献。但是，当吹美一些国家从 17 世纪中叶开始确立资本主义生产方


式，从18世纪60年代开始工业革命的时候，中国最后一个封建王朝——清朝的统治者却夜郎自大、自我封闭，拒绝扩大与外国交往。原来文明程度落后于中国的欧美国家，这时跑到了中国前面。也就是说，导致中国落后于时代的原因是中国封建统治者夜郎自大、闭关锁国，A正确。落后就要挨打。1840年，急于向外扩张的英国发动侵略中国的鸦片战争，中国历史的发展从此发生重大转折。B、C、D都不是导致中国落后于时代的原因，而是后果，故错误。


\subsection*{10. 【答案】D}

\subsection*{【考点知识还原】}

见表 82


\subsection*{【考点点拨】}

八七会议


\subsection*{【试题简析】}

大革命失败以后，全国完罩在白色恐怖之下，中国共产党及其领导的革命运动遭到严酷镇压。在革命的危急关头，1927年7月中旬，中共中央政治局临时常委会决定了三件大事：将党所掌握和影响的部队向南昌集中，准备起义；组织湘、鄂、赣、粵四省的农民发动秋收起义；召集中央紧急会议，讨论和决定大革命失败后的新方针。C错误。同年8月7日，中共中央在汉口秘密召开紧急会议（八七会议），会议着重批评了大革命后期党内以陈独秀为代表的右倾机会主义错误，确定了土地革命和武装起义的方针，这是一个正确的方针，是党在付出血的代价后换得的正确结论。出席这次会议的毛泽东在发言中突出地强调：“以后要非常注意军事。须知政权是由枪杆子中取得的。”会议还提出了“找着新的道路”的任务。八七会议给正处在思想混乱和组织涣散中的中国共产党指明了新的出路，这是由大革命失败到土地革命战争兴起的历史性转变。A错误，D正确。八七会议实际上还是坚持城市中心的思想，并没有阐明党的工作重心从城市转移到农村的重要性。八七会议后党领导举行的秋收起义、广州起义和其他许多地区起义大多数失败了的事实证明，在当时的客观条件下，中国共产党人不可能像俄国十月革命那样通过首先占领中心城市来取得革命在全国的胜利，党迫切需要找到适合中国国情的革命道路。毛泽东同志领导军民在井冈山建立第一个农村革命根据地，党领导人民打土豪、分田地。从进攻大城市转为向农村进军，是中国革命具有决定意义的新起点。B错误。


\subsection*{11. 【答案】A}

\subsection*{【考点知识还原】}

见表87


\subsection*{【考点点拨】}

延安精神


\subsection*{【试题简析】}

2022年10月27日，党的二十大闭幕不到一周，习近平带领中共中央政洽局常委专程从北京前往陕西延安，瞻仰延安革命纪念地，重温革命战争时期党中央在延安的峥嵘岁月，缅怀老一辈革命家的丰功伟绩，宣示新一届中央领导集体赓续红色血脉、传承奋斗精神，在新的赶考之路上向历史和人民交出新的优异答卷的坚定信念。习近平指出，延安是中国革命的圣地、新中国的摇篮。从 1935年到1948年，党中央和毛泽东等老一辈革命家在延安生活和战斗了13年，领导中国革命事业从低潮走向高潮、实现历史性转折，扭转了中国前途命运。习近平强调，在延安时期形成和发扬的光荣传统和优良作风，培育形成的以坚定正确的政治方向、解放思想实事求是的思想路线、全心全意为人民服务的根本宗旨、自力更生艰苦奋斗的创业精神为主要内容的延安精神，是党的宝贵精神财富，要代代传承下去。延安时期，党提出全心全意为人民服务的根本宗旨并写入党章，强调共产党“这个队伍完全是为着解放人民的，是彻底地为人民的利益工作的”，要求党的干部“把屁股端端地坐在老百姓的这一面”，形成了“只见公仆不见官”的生动局面。全党同志要站稳人民立场，践行党的宗旨，贯彻党的群众路线，保持党同人民群众的血肉联系，自觉把以人民为中心的发展思想贯穿到各项工作之中，扎实推进共同富裕，让现代化建设成果更多更公平惠及全体人民。A 正确。


\subsection*{12. 【答案】B}

\subsection*{【考点知识还原】}

见表92


\subsection*{【考点点拨】}

中国进入社会主义社会最主要的标志


\subsection*{【试题简析】}

划分社会形态的依据是社会的经济基础特别是生产关系的性质。1956年年底，随着社会主义改造的完成，以生产资料公有制、按劳分配和计划经济体制为特征的社会主义经济制度建立起来，这是中国进入社会主义社会最主要的标志。B正确。人民民主专政国家政权的建立，是在1949年10月；民主革命遗留任务的完成是在1952年；社会主义基本政治制度中的中国共产党领导的多党合作和政治协商制度、民族区域自治制度是在1949年第一届政治协商会议的召开及其制定的《中国人民政治协商会议共同纲领》中确立的，人民代表大会制度是在1954年9月确立的，A、C、D错误。


\subsection*{13. 【答案】C}

\subsection*{【考点知识还原】}

见表103


\subsection*{【考点点拨】}

新时代公民道德建设的着力点


\subsection*{【试题简析】}

2001年，党中央颁布《公民道德建设实施纲要》，对在社会主义市场经济条件下加强公民道德建设提供了重要指导，有力促进了社会主义精神文明建设。2019年中共中央、国务院颁布了《新时代公民道德建设实施纲要》，要求各地区、各部门结合实际认真贯彻落实。《新时代公民道德建设实施纲要》强调，中国特色社会主义进入新时代，加强公民道德建设、提高全社会道德水平，是全面建成小康社会、全面建设社会主义现代化强国的战略任务，是适应社会主要矛盾变化、满足人民对美好生活向在的迫切需要，是促进社会全面进步、人的全面发展的必然要求。纲要明确指出，新时代公民道德建设，要把社会公德、职业道德、家庭美德、个人品德建设作为着力点。推动践行以文明礼貌、助人为乐、爱护公物、保护环境、遵纪守法为主要内容的社会公德，鼓励人们在社会上做一个好公民；推动践行以爱岗敬业、诚实守信、办事公道、热情服务、奉献社会为主要内容的职业道德，鼓励人们在工作中做一个好建设者；推动践行以尊老爱幼、男女平等、夫妻和睦、勤俭持家、邻里互助为主要内容的家庭美德，鼓励人们在家庭里做一个好成员；推动践行以爱国奉献、明礼遵规、勤劳善良、宽厚正直、自强自律为主要内容的个人品德，鼓励人们在日常生活中养成好品行。C正确。


\subsection*{14. 【答案】A}

\subsection*{【考点知识还原】}

见表106


\subsection*{【考点点拨】}

坚持依法治国首先要坚持依宪治国


\subsection*{【试题简析】}

党的二十大报告指出，全面依法治国是国家治理的一场深刻革命，关系党执政兴国，关系人民幸福安康，关系党和国家长治久安。必须更好发挥法治固根本、稳预期、利长远的保障作用，在法治轨道上全面建设社会主义现代化国家。我们要坚持走中国特色社会主义法治道路，建设中国特色社会主义法治体系、建设社会主义法治国家，围绕保障和促进社会公平正义，坚持依法治国、依法执政、依法行政共同推进，坚持法治国家、法治政府、法治社会一体建设，全面推进科学立法、严格执法、公正司法、全民守法，全面推进国家各方面工作法治化。要完善以宪法为核心的中国特色社会主义法律体系。坚持依法治国首先要坚持依宪治国，坚持依法执政首先要坚持依宪执政，坚持宪法确定的中国共产党领导地位不动摇，坚持宪法确定的人民民主专政的国体和人民代表大会制度的政体不动摇。A正确。*15.【答案】B


\subsection*{15. 【答案】B}

\subsection*{【考点点拨】}

“形势与政策”


\subsection*{【试题简析】}

北京冬奥会、冬残奥会是新冠病毒感染疫情发生以来首次如期举办的全球综合性


体育盛会，是一场和平友谊的盛会、一场团结合作的盛会、一场鼓舞世界的盛会。北京冬奥会、冬残奥会的成功举办，促进了不同文明交流互鉴，为推动全球团结合作、共克时艰发挥了重要作用，也为动荡不安的世界带来了信心和希望，向世界发出了“一起向未来”的时代强音！2022年4月8日，北京冬奥会、冬残奥会总结表彰大会在北京举行。习近平强调，伟大的事业孕育伟大的精神，伟大的精神推进伟大的事业。北京冬奥会、冬残奥会广大参与者珍惜伟大时代赋予的机遇，在冬奥申办、筹办、举办的过程中，共同创造了北京冬奥精神。北京冬奥精神就是胸怀大局、自信开放、迎难而上、追求卓越、共创未来。B正确。“更快、更高、更强、更团结”是奥林匹克格言；“勇气、决心、激励、平等”是残疾人奥林匹克（残奥）价值观；“一起向未来”是2022年北京冬奥会、冬残奥会的主题口号。A、C、D错误。*16.【答案】D


\subsection*{16. 【答案】D}

\subsection*{【考点点拨】}

“形势与政策”


\subsection*{【试题简析】}

2022年11月11日，第二十五次中国一东盟（10+1）领导人会议在東埔寨首都金边举行。会议通过了《关于加强中国一东盟共同的可持续发展联合声明》《纪念〈南海各方行为宣言〉签署二十周年联合声明》《中国一东盟粮食安全合作的联合声明》等成果文件，欢迎通过《中国一东盟全面战略伙伴关系行动计划》，宣布2023年为“中国东盟农业发展和粮食安全合作年”。D正确，C错误。《落实〈南海各方行为宣言〉指导方针》与《中国和东盟国家外交部长关于全面有效落实〈南海各方行为宣言〉联合声明》分别于2011年、2016年在中国一东盟（10+1）领导人会议上通过，A、B错误。


\subsection*{二、多项选择题}

\subsection*{17. 【答案】ABC}

\subsection*{【考点知识还原】}

见表 19


\subsection*{【考点点拨】}

无产阶级领袖在历史上的作用


\subsection*{【试题简析】}

个人在历史上的作用存在差别。有的人作用大些，可称为“历史人物”；有的人作用小些，可称为“普通个人”。历史人物是一定历史事件的主要倡导者、组织领导者或思想理论、科学文化的重要代表人物。历史人物对历史发展有深刻影响，甚至有时能够决定个别历史事件的结局，从而导致历史发生这样或那样的重大变化。从其发挥作用的性质来看，有些历史人物起推动历史前进的进步作用，有些起阻碍历史前进的反动作用。杰出人物是历史人物中对推动历史发展作出重要贡献或起重要作用的人。在历史发展进程中，新的历史任务往往是由具有进步意义的历史人物首先发现或提出来的。他们比一般人站得高、看得远，解决历史任务的愿望比一般人强烈。先进阶级的政治代表人物，特别是无产阶级的领袖人物，所提出的思想能够成为社会变革的先导。他们为群众指明革命斗争的方向，在革命斗争中起着领导核心的作用。有些占统治地位的剥削阶级的政治代表，在特定的条件下运用其权力满足了社会某些方面的需要，对历史发展也会起到某种作用甚至是重大的促进作用。A、B、C正确。但是必须明确，不管什么样的历史人物，在历史上发挥什么样的作用，都要受到社会发展客观规律的制约，而不能决定和改变历史发展的总进程和总方向。杰出人物会因其智慧、性格因素对社会进程产生影响，但这些作用仅仅是历史进程中的偶然现象，只能成为社会发展的个别原因。他们凭借自己的才能，虽然也能使具体历史事件的外貌或某些后果发生变化，但终究不能改变历史发展的基本方向。无产阶级领袖不同于以往历史上的杰出人物，他们所代表的是历史上最革命、最先进的阶级，在革命和建设中发挥了重大作用。但是，他们也不能创造人类社会的历史。创造人类社会历史的是人民。毛泽东说：“人民，只有人民，才是创造世界历史的动力。”D错误。


\subsection*{18. 【答案】BD}

\subsection*{【考点知识还原】}

见表13


\subsection*{【考点点拨】}

价值评价的特点


\subsection*{【试题简析】}

价值评价是主体对客体的价值以及价值大小所作的评判或判断，因而也被称作价值判断。价值评价是对客观价值关系的主观反映。在价值评价活动中，主体以自身的需要作为内在尺度运用于所评价的客体，因此，评价的差异性是无法避免的。能否作出正确判断，取决于人对客体和主体的双重认识，这种认识就包括对客体属性和规律的认识，A错误。价值评价的主要特点是：第一，评价以主客体之间的价值关系为认识对象。价值评价以客体和主体之间的价值关系为反映内容，主体的意向、愿望和要求包含在其中，追求“应该怎样”，以求“善”和“美”为认识目的。第二，评价结果与评价主体直接相关，受主体意志的影响。B正确。第三，评价结果的正确与否依赖于对客体状况和主体需要的认识。价值关系中的主体是在一定社会关系中从事实践的具体的人，随着实践和历史的发展，主体和客体以及主客体之间的关系发生变化，导致人们对客体价值的判断也发生改变。D正确。第四，价值评价有科学与非科学之别。评价结果有正确与错误之分，价值评价也有科学与非科学之别。C错误。


\subsection*{19. 【答案】ABC}

\subsection*{【考点知识还原】}

见表28


\subsection*{【考点点拨】}

平均利润


\subsection*{【试题简析】}

资本主义生产的目的是获得利润。为了得到尽可能高的利润率和尽可能多的利润，不同生产部门的资本家之间必然展开激烈的竞争，大量资本必然从利润率低的部门转投到利润率高的部门，从而导致利润率平均化。在利润率平均化的过程中，形成了社会的平均利润率，按照平均利润率计算和获得的利润，叫作平均利润。利润转化为平均利润，是剩余价值规律和竟争规律作用的必然结果，体现着不同部门的资本家集团要求按照等量资本获得等量利润的原则来瓜分剩余价值的关系。在利润率平均化的过程中，产业资本家得到产业利润，商业资本家得到商业利润，银行资本家得到银行利润，农业资本家得到农业利润，这些不同部门的资本家瓜分到的利润只是平均利润。平均利润率规律的作用表明，平均利润率是剩余价值总量与社会总资本的比率。每个资本家所得利润多少不仅取决于他对本企业工人的剥削程度，而且还取决于整个资产阶级对整个工人阶级的剥削程度。资本家之间在瓜分剩余价值上固然有一定程度的利害冲突，但在加强对工人阶级的剥削以榨取更大量的剩余价值这一点上，有着共同的阶级利益。A、B、C正确，D错误。


\subsection*{【名师点拨】}

本题的知识点在 2019 年第19题考查过一次。


\subsection*{20. 【答案】AC}

\subsection*{【考点知识还原】}

见表30


\subsection*{【考点点拨】}

金融垄断资本发展


\subsection*{【试题简析】}

第二次世界大战后，在美国的主导下建立了国际货币体系，即布雷顿森林体系，该体系对促进世界经济的恢复和发展发挥了重要作用，但也维护了美国的世界经济霸权地位。20世纪 70年代初，由于资本主义发展不平衡的加深和国际货币体系内在矛盾的激化，布雷顿森林体系崩溃。随后，西方国家普遍走上金融自由化和金融创新的道路。金融自由化与金融创新是金融垄断资本得以形成和壮大的重要制度条件。金融自由化的主要内容是：各国政府放松对银行利率的管制，实行浮动汇率制度，取消外汇管制，金融市场相互开放。金融创新的主要内容是：包括远期合约、期货合约、期权合约、掉期合约等在内的信用风险防范工具和融资技术不断推陈出新，金融机构开始突破原有的专业分工界限，综合经营各种金融业务，金融工具不断创新，传统信贷业务逐年减少，债券业务迅速增长，融资方式证劵化趋势迅猛发展。A、C正确。


\subsection*{【名师点拨】}

本题的知识点在2016 年第4题考查过一次。


\subsection*{21. 【答案】AB}

\subsection*{【考点知识还原】}

见表36


\subsection*{【考点点拨】}

战时共产主义政策


\subsection*{【试题简析】}

苏维埃政权的建立和现固，引起了国内外敌对势力的恐惧、仇视和反抗。从1918年下半年起，西方列强纠合14个国家发动了对苏维埃政权的武装干涉，同时俄国爆发了国内战争。为了打赢这场战争，捍卫年轻的苏维埃政权，建设社会主义新社会，从1918年夏到1921年春，列宁领导的苏维埃政权实行了以余粮收集制和取消商品货币关系为主要特征的战时共产主义政策，在经济上采取一系列特殊的又带有强制性的非常措施。战时共产主义政策是苏维埃俄国面临帝国主义武装干涉和国内战争而被迫采取的一种临时性政策。正如列宁后来所说的：“为了拯救国家，拯救军队，拯救工农政权，当时必须这样做。”战时共产主义政策的实施，对于粉碎国际帝国主义的武装干涉和国内反革命叛乱，保卫新生的苏维埃政权，发挥了重要作用。但是，在1921年初转人和平经济建设时期后，如不改变政策、继续实行战时共产主义政策，将会失去广大群众，丢掉阶级基础。1921年3月，俄共（布）召开十大，毅然决定从战时共产主义政策过渡到实行以发展商品经济为主要特征的新经济政策。A、B正确。


\subsection*{22. 【答案】ABC}

\subsection*{【考点知识还原】}

见表39


\subsection*{【考点点拨】}

毛泽东对马克思主义中国化时代化的主要贡献


\subsection*{【试题简析】}

毛泽东是马克思主义中国化时代化的伟大开拓者，是毛泽东思想的主要创立者。在中国共产党历史上，毛泽东第一个明确提出了“马克思主义中国化”的科学命题和重大任务，深刻论证了马克思主义中国化的必要性和极端重要性，系统阐述了马克思主义中国化的科学内涵和实现马克思主义中国化的正确途径，开辟了马克思主义在中国发展的宽广道路，为党领导的革命和建设事业的发展奠定了坚实的思想理论基础。毛泽东为实现马克思主义中国化时代化的历史任务进行了艰苦的探索，使马克思主义在中国生根、开花、结果。A、B、C正确。中国共产党的历史，是一部马克思主义中国化时代化的历史，是一部不断推进理论创新、进行理论创造的历史。在这个过程中，不可避免地会出现各种各样错误的认识，毛泽东和毛泽东思想的创立，不可能全面清除马克思主义中国化认识上的各种障碍，D错误。


\subsection*{23. 【答案】ABCD}

\subsection*{【考点知识还原】}

见表55


\subsection*{【考点点拨】}

历史地、发展地、辩证地认识和把握我国社会存在的各类资本及其作用


\subsection*{【试题简析】}

2022年4月29日下午，中共中央政治局就依法规范和引导我国资本健康发展进行第38次集体学习。中共中央总书记习近平在主持学习时发表了重要讲话。习近平指出，要历史地、发展地、辩证地认识和把握我国社会存在的各类资本及其作用。在社会主义市场经济体制下，资本是带动各类生产要素集聚配置的重要纽带，是促进社会生产力发展的重要力量，要发挥资本促进社会生产力发展的积极作用。同时，必须认识到，资本具有逐利本性，如不加以规范和约束，就会给经济社会发展带来不可估量的危害。我们要立足新发展阶段、贯彻新发展理念、构建新发展格局、推动高质量发展，正确处理不同形态资本之间的关系，在性质上要区分，在定位上要明确，规范和引导各类资本健康发展。习近平强调，要正确处理资本和利益分配问题。我国社会主义的国家性质决定了我们必须坚持按劳分配为主体、多种分配方式并存，在社会分配中体现人民至上。要注重经济发展的普惠性和初次分配的公平性，既注重保障资本参与社会分配获得增殖和发展，更注重维护按劳分配的主体地位，坚持发展为了人民、发展依靠人民、发展成果由人民共享，坚定不移走全体人民共同富裕的道路。习近平指出，要深化资本市场改革。要继续完善我国资本市场基础制度，更好发挥资本市场功能，为各类资本发展释放出更大空间。A、B、C、D全选。


\subsection*{24. 【答案】BCD}

\subsection*{【考点知识还原】}

见表68


\subsection*{【考点点拨】}

中国国际进口博览会作用


\subsection*{【试题简析】}

2022年11月4日晚，国家主席习近平以视频方式出席在上海举行的第五届中国国际进口博览会开幕式并发表题为《共创开放繁荣的美好未来》的致辞。习近平指出：“5年前，我宣布举办进博会，就是要扩大开放，让中国大市场成为世界大机遇。现在，进博会已经成为中国构建新发展格局的窗口、推动高水平开放的平台、全球共享的国际公共产品。”习近平强调，开放是人类文明进步的重要动力，是世界繁荣发展的必由之路。当前，世界百年未有之大变局加速演进，世界经济复苏动力不足。我们要以开放纾发展之困、以开放汇合作之力、以开放聚创新之势、以开放谋共享之福，推动经济全球化不断向前，增强各国发展动能，让发展成果更多、更公平惠及各国人民。B、C、D正确，A错误。


\subsection*{25. 【答案】ABCD}

\subsection*{【考点知识还原】}

见表64


\subsection*{【考点点拨】}

湿地保护与生态文明建设


\subsection*{【试题简析】}

2022年11月5日，习近平以视频方式出席在武汉举行的《湿地公约》第十四届缔约方大会开幕式并发表致辞。习近平指出：“古往今来，人类逐水而居，文明伴水而生，人类生产生活同湿地有着密切联系。本次大会以‘珍爱湿地，人与自然和谐共生’为主题，共谋湿地保护发展，具有十分重要的意义。我们要深化认识、加强合作，共同推进湿地保护全球行动。”习近平提出要凝聚珍爱湿地全球共识，深怀对自然的敬畏之心，减少人类活动的干扰破坏，守住湿地生态安全边界，为子孙后代留下大美湿地；要推进湿地保护全球进程，加强原真性和完整性保护，把更多重要湿地纳入自然保护地，健全合作机制平台，扩大国际重要湿地规模；要增进湿地惠民全球福祉，发挥湿地功能，推进可持续发展，应对气候变化，保护生物多样性，给各国人民带来更多实惠。A、B、C、D全选。


\subsection*{26. 【答案】BD}

\subsection*{【考点知识还原】}

见表38


\subsection*{【考点点拨】}

坚持和发展马克思主义必须坚持“两个结合”


\subsection*{【试题简析】}

党的二十大报告指出，中国共产党人深刻认识到，只有把马克思主义基本原理同中国具体实际相结合、同中华优秀传统文化相结合，坚持运用辩证唯物主义和历史唯物主义，才能正确回答时代和实践提出的重大问题，才能始终保持马克思主义的蓬勃生机和旺盛活力。坚持和发展马克思主义，必须同中华优秀传统文化相结合。只有植根本国、本民族历史文化沃土，马克思主义真理之树才能根深叶茂。中华优秀传统文化源远流长、博大精深，是中华文明的智慧结晶，其中蕴含的天下为公、民为邦本、为政以德、革故鼎新、任人唯贤、天人合


\subsection*{一、自强不息、厚德载物、讲信修睦、亲仁善邻等，是中国人民在长期生产生活中积累字宙}

观、天下观、社会观、道德观的重要体现，同科学社会主义价值观主张具有高度契合性。我们必须坚定历史自信、文化自信，坚持古为今用、推陈出新，把马克思主义思想精髓同中华优秀传统文化精华贯通起来、同人民群众日用而不觉的共同价值观念融通起来，不断赋予科学理论鲜明的中国特色，不断夯实马克思主义中国化时代化的历史基础和群众基础，让马克思主义在中国牢牢扎根。B、D正确。中华优秀传统文化源远流长，已经有几千年的历史，马克思主义诞生至今才170多年，所以，中华优秀传统文化的这些思想观念，不可能和马克思主义理论一脉相承；这些思想观念同科学社会主义价值观主张具有高度契合性，不是、也不可能完整体现马克思主义的立场、观点和方法，A、C错误。


\subsection*{27. 【答案】ACD}

\subsection*{【考点知识还原】}

见表78


\subsection*{【考点点拨】}

革命派与改艮派的论战


\subsection*{【试题简析】}

在资产阶级民主革命思潮广泛传播、革命形势日益成熟的时候，康有为、梁启超等人坚持走改良道路，反对用革命手段推翻清朝统治。1905年至1907年间，围绕中国究竟是采用革命手段还是改良方式这个问题，革命派与改良派分别以《民报》《新民丛报》为主要舆论阵地，展开了一场大论战。投入这场论战的还有其他十几种报刊。双方论战问题包括要不要以革命手段推翻清王朝，要不要推翻帝制、实行共和，要不要进行社会革命，焦点是要不要以革命手段推翻清王朝，A正确。这场论战具有重大的意义。通过这场论战，划清了革命与改良的界限，传播了民主革命思想，促进了革命形势的发展。但这场论战也暴露了革命派在思想理论方面的弱点。比如，他们主张推翻清政府，但对“革命是否会招致帝国主义干涉”的问题不敢作出理直气壮的正面回答，只是希望通过“有秩序的革命”来避免动乱和帝国主义的干涉。他们所说的“国民”，主要还是指资产阶级及其知识分子，而不是广大的劳动群众。他们对封建地主土地所有制是否应该改革的问题也是语焉不详，并且反对贫苦农民“夺富人之田为已有”。这些理论和认识的局限不可避免地会影响辛亥革命的进程和结局。C、D正确。这场论战是资产阶级内部革命派与改良派的论战，不是资产阶级思想与封建主义思想的交锋。“百日维新”前的维新派与守旧派的论战，实质上是资产阶级思想与封建主义思想在中国的第一次正面交锋。B错误。


\subsection*{28. 【答案】CD}

\subsection*{【考点知识还原】}

见表83


\subsection*{【考点点拨】}

毛泽东的著作——《反对本本主义》


\subsection*{【试题简析】}

1930年5月，毛泽东在《反对本本主义》一文中，阐明了坚持辩证唯物主义的思想路线即坚持理论与实际相结合的原则的极端重要性，提出了 “没有调查，没有发言权”和“中国革命斗争的胜利要靠中国同志了解中国情况”的重要思想，表现了毛泽东开辟新道路、创造新理论的革命首创精神。C、D 正确。“关心群众生活，注意工作方法”，是毛泽东1934年1月27日在江西瑞金召开的第二次全国工农兵代表大会上所作的结论的一部分。“要注意大城市和交通要道的工作，要把城市工作和根据地工作提到同等重要的地位”，是1944年5月20日毛泽东在中央党校作时局问题报告时提出的。1944年6月5日，中共中央发布了毛泽东起草的《中共中央关于城市工作的指示》，实际上是对毛泽东提出的“要注意大城市和交通要道的工作，要把城市工作和根据地工作提到同等重要的地位”思想的具体化。A、B错误。


\subsection*{29. 【答案】ABCD}

\subsection*{【考点知识还原】}

见表94


\subsection*{【考点点拨】}

中共十一届三中全会作出实行改革开放的历史性决策的依据


\subsection*{【试题简析】}

2018年12月18日，习近平在庆祝改革开放40周年大会上的讲话中指出，党的十一届三中全会是在党和国家面临何去何从的重大历史关头召开的。在邓小平同志领导下和老一辈革命家支持下，党的十一届三中全会冲破长期“左”的错误的严重束缚，批评“两个凡是”的错误方针，充分肯定必须完整、准确地掌握毛泽东思想的科学体系，高度评价关于真理标准问题的讨论，果断结束“以阶级斗争为纲”口号，重新确立马克思主义的思想路线、政治路线、组织路线。从此，我国改革开放拉开了大幕。习近平在讲话中强调，我们党作出实行改革开放历史性决策，是基于对党和国家前途命运的深刻把握，是基于对社会主义革命和建设实践的深刻总结，是基于对时代潮流的深刻洞察，是基于对人民群众期盼和需要的深刻体悟。改革开放是我们党的一次伟大觉醒，正是这个伟大觉醒孕育了我们党从理论到实践的伟大创造。改革开放是中国人民和中华民族发展史上一次伟大革命，正是这个伟大革命推动了中国特色社会主义事业的伟大飞跃！A、B、C、D全选。


\subsection*{30. 【答案】ABD}

\subsection*{【考点知识还原】}

见表103


\subsection*{【考点点拨】}

社会主义道德先进性的体现


\subsection*{【试题简析】}

道德同其他社会意识形态一样，不是亘古不变的。迄今为止，人类社会先后经历了五种基本社会形态，与此相适应，出现了原始社会的道德、奴隶社会的道德、封建社会的道德、资本主义社会的道德、社会主义社会的道德。在社会主义社会，有一部分先进分子还身体力行共产主义道德。与以往社会的道德形态相比，社会主义道德具有显著的先进性特征。这种先进性主要体现在以下几个方面。首先，社会主义道德是社会主义经济基础的反映。在以生产资料公有制为主体的社会主义社会，广大人民不仅在政治上实现了当家作主，而且在道德上实现了由被动到主动的转变。其次，社会主义道德是对人类优秀道德资源的批判继承和创新发展。以当代中国的社会主义道德体系为例，我们今天倡导的社会主义道德规范，不仅与中华传统美德相承接，与中国共产党人在革命战争年代创立的革命道德相延续，同时也是对人类优秀道德成果的吸收和借鉴。最后，社会主义道德克服了以往阶级社会道德的片面性和局限性，坚持以为人民服务为核心，坚持以集体主义为原则，展现出真实而强大的道义力量。A、B、D正确。调节社会一切行为规范的准则，不仅有道德，还有法律，C错误。


\subsection*{31. 【答案】AC}

\subsection*{【考点知识还原】}

见表 109


\subsection*{【考点点拨】}

法律义务的特点


\subsection*{【试题简析】}

法律义务与法律权利相对应。法律义务是指由一定的社会物质生活条件所制约的社会责任，是保证法律所规定的义务人按照权利人要求从事一定行为或不从事一定行为以满足权利人利益的法律手段。只有承担法律义务的人履行法律义务，享有法律权利的人才能实现自己的合法权益。法律义务具有法定的强制性，违反法律义务必须承担法律责任。D错误。法律义务具有以下4个特点。第一，法律义务是历史的。法律义务的内容和履行方式随着经济社会的发展和人权保障的进步而不断调整和变化。第二，法律义务源于现实需要。一个国家或地区的制度性质、历史传统、文化背景、宗教信仰和安全形势等因素，会对法律义务的设定产生重要影响。第三，法律义务必须依法设定。法律义务必须由具有法定职权的国家机关依照法律程序设定，其他国家机关不得对公民违法设定法律义务。坚持义务法定，是建设法治国家和保障人权的重要方面。第四，法律义务可能发生变化。公民和社会组织承担的法律义务，在履行的过程中可能会基于法定情形而变更、消灭，或产生新的法律义务。A、C正确，B错误。*32.【答案】ABCD


\subsection*{32. 【答案】ABCD}

\subsection*{【考点点拨】}

“形势与政策”


\subsection*{【试题简析】}

2022年8月10日，国务院台湾事务办公室、国务院新闻办公室发表《台湾问题与新时代中国统一事业》白皮书，进一步重申台湾是中国的一部分的事实和现状，展现中国共产党和中国人民追求祖国统一的坚定意志和坚强决心，阐述中国共产党和中国政府在新时


代推进实现祖国统一的立场和政策。在当前错综复杂的国际形势和台海形势下，发表该白皮书，有利于揭批“台独”分裂势力和外部势力勾连挑衅，企图损害中国主权和领土完整、阻挠破坏中国统一进程的恶劣言行，揭露他们的政治本质和险恶用心，体现中国人民维护国家主权和领土完整的决心不可动摇、意志坚如磐石；有利于展现中国共产党和中国政府愿继续以最大诚意、尽最大努力争取和平统一的立场和态度，最大限度地争取两岸同胞特别是台湾同胞与国际社会的理解和支持；有利于提振全党全国各族人民矢志追求国家统一的精气神，增强岛内和海外反“独”促统力量的信心和勇气，凝聚支持和促进祖国统一的磅礴伟力。A、B、C、D全选。


\subsection*{33. 【答案】BCD}

\subsection*{【考点点拨】}

“形势与政策”


\subsection*{【试题简析】}

2022年11月15日，二十国集团领导人第十七次峰会在印度尼西亚巴厘岛举行。国家主席习近平出席并发表题为《共迎时代挑战 共建美好未来》的重要讲话。这是中共二十大胜利召开后中国最高领导人首次出访，是新冠病毒感染疫情近3年来习近平首次线下出席多边会议，也是世界又一次站在十字路口的关键时刻中国开展的重大外交行动。习近平在《共迎时代挑战 共建美好未来》的重要讲话中强调，二十国集团成员都是世界和地区大国，应该体现大国担当，发挥表率作用，为各国谋发展，为人类谋福祉，为世界谋进步。为此，习近平提出3点倡议：推动更加包容的全球发展，推动更加普惠的全球发展，推动更有韧性的全球发展。B、C、D正确。


\subsection*{三、材料分析题}

\subsection*{34. 【答案】}

（1）科学家在实践中获得对盐碱地的土壤结构、不同区域盐碱特征等感性认识，运用科学思维方法，将这些感性认识上升为治理盐碱地的治理原则和治理方案，实现从感性认识到理性认识的飞跃。科学家将盐碱地治理原则和治理方案用以指导实践，建立起不同区域生态修复盐碱地系统工程技术体系，实现从理性认识到实践的飞跃。科学家在解决盐碱地治理的问题上，反复实验验证，不断优化治理方案，深化认识，体现了认识与实践矛盾运动的反复性和上升性。


（2）万事万物是相互联系、相互依存的，只有用普遍联系的、全面系统的、发展变化的观点观察事物，才能把握事物发展规律。系统观念是运用事物普遍联系观点，把握系统与要素、要素与要素、结构与层次、系统与环境之间相互联系和作用的动态过程，是具有基础性的思想和工作方法。盐碱荒地是众多因素相互联系、相互影响的系统整体，不同盐碱区域土壤结构和盐碱化程度差异大，盐碱地改良必须坚持系统观念，协同工程设施与土壤改良方法，系统优化治理方案，实现治理成效最优化。


\subsection*{【考点点拨】}

实践与认识的辩证运动，认识的反复性和无限性，事物的普遍联系，坚持系统观念


\subsection*{【试题简析】}

第一问应该从实践与认识的辩证运动这个原理所包含的两个角度作答。先是阐述人们认识具体事物的辩证运动过程，即从实践到认识，再从认识到实践，完成了两次飞跃；然后再指出针对实践和认识运动过程的向前推移、向前发展而言，认识运动并没有完成，它是一个反复循环和无限发展的过程。具体组织答案时要紧紧抓住以下两个踩分点：一是认识过程的两次飞跃。科学家在实践中获得对盐碱地的土壤结构、不同区域盐碱特征等感性认识，运用科学思维方法，将这些感性认识上升为洽理盐碱地的洽理原则和洽理方案，实现从感性认识到理性认识的飞跃。科学家将盐碱地治理原则和治理方案用以指导实践，建立起不同区域的生态修复盐碱地系统工程技术体系，实现从理性认识到实践的飞跃。二是认识的反复性和无限性。科学家在解决盐碱地洽理的问题上，反复实验验证，不断优化治理方案，深化认识，体现了实践与


认识矛盾运动的反复性和上升性。第二问与第一问的问法略有不同。第一问要求写出材料所体现的认识论原理（“是什么”），因此应该在结合题目语料的同时回答相应的哲学原理，即将原理和材料整合在一起作答。而第二问是要回答“为什么”，因此更好的做法是先回答哲学原理本身，再结合材料简要说明。按照这个答题思路，第二问要从两个方面回答：首先，要说明系统观念哲学依据（事物普遍联系原理）及其含义，指出万事万物是相互联系、相互依存的，只有用普遍联系、全面系统的、发展变化的观点观察事物，才能把握事物发展规律。系统观念是运用事物普遍联系观点，把握系统与要素、要素与要素、结构与层次、系统与环境之间相互联系和作用的动态过程，是具有基础性的思想和工作方法。然后再结合材料，指出盐碱荒地是众多因素相互联系、相互影响的系统整体，不同盐碱区域土壤结构和盐碱化程度差异大，盐碱地改良必须坚持系统观念，协同工程设施与土壤改良方法，系统优化治理方案，实现治理成效最优化。第二问踩分点有3个：原理——事物普遍联系原理；概念——系统思维能力；系统思维在盐碱地治理过程中的体现及效果。


\subsection*{【名师点拨】}

2023年的这道马原哲字题应当讲是比较容易的。一是本题的两问都明确了要考查的原理；二是我时2023版预测春把认识论和系统观念作为重中之重。考生作答时需要注意两点：一是认真审题，如第一问考的是实践与认识的辩证运动，因此只回答二者的辩证关系是不够的，必须要阐述认识辩证运动发展的过程；二是要将具体的题目语料带人到哲学原理中，结合题目来作答，比如 【答案要点】第一问的第一句就点明了是关于什么的感性认识（对盐碱地的土壤结构、不同区域盐碱特征等），形成了怎样的理性认识（治理盐碱地的治理原则和治理方案）。


\subsection*{35. 【答案】}

（1）人民性是马克思主义的本质属性，坚持人民至上是中国共产党人的根本政治立场。坚持以人民为中心的发展思想，既是我们党领导现代化建设的出发点和落脚点，也是新发展理念的“根”和“魂”。以人民为中心的发展思想体现在经济社会发展各个环节。十年的发展成就表明，我们党坚持发展为了人民、发展依靠人民、发展成果由人民共享，让现代化建设成果更多更公平惠及全体人民；坚持在发展中保障和改善民生，着力解决好人民群众急难愁盼问题，使人民群众获得感、幸福感、安全感更加充实；坚持相信和依靠人民群众、尊重群众的历史首创精神，从群众中汲取无穷 的智慧和力量。


（2） 新时代十年来，党对中国式现代化的深远意义、科学内涵、本质要求及目标蓝图、发展规划、战略导向等方面都作了丰富和完善，成功推进和拓展了中国式现代化。党和国家事业取得历史性成就、发生历史性变革，为全面建成社会主义现代化国家奠定了坚实基础和创造了有利条件；创立习近平新时代中国特色社会主义思想，为全面建成社会主义现代化国家提供根本遵循；提出立足新发展阶段、贯彻新发展理念、构建新发展格局、推进高质量发展；作出全面建成社会主义现代化强国新的两步走战略安排，把第二个百年奋斗目标充实提升为建成富强民主文明和谐美丽的社会主义现代化强国。


\subsection*{【考点点拨】}

以人民为中心的发展思想，党的十八大以来所理论和实践创新突破成功推进和拓展了中国式现代化


\subsection*{【试题简析】}

本题考查的也是考生比较熟悉的知识点。第一问考查以人民为中心的发展思想。回答这个问题，一定要对提问的要求作必要的延伸，就是说，这个提问不能只回答新时代十年我们是如何体现以人民为中心的发展思想的，还要回答为什么要坚持以人民为中心的发展思想以


及什么是以人民为中心的发展思想。搞清楚这一点是很重要的。所以，回答第一问，首先要说明以人民为中心的发展思想的重要性及其含义，指出人民性是马克思主义的本质属性，坚持人民至上是中国共产党人的根本政治立场。坚持以人民为中心的发展思想，既是我们党领导现代化建设的出发点和落脚点，也是新发展理念“根”和“魂”。所谓以人民为中心的发展思想，就是坚持发展为了人民，发展依靠人民，发展成果由人民共享。然后，再结合已经掌握新发展理念和材料1提供的信息，说明新时代十年“党和国家事业取得历史性成就、发生历史性变革”所体现以人民为中心的发展思想，即十年的发展成就表明，我们党让现代化建设成果更多更公平惠及全体人民；坚持在发展中保障和改善民生，着力解决好人民群众急难愁盼问题，使人民群众获得感、幸福感、安全感更加充实；坚持相信和依靠人民群众、尊重群众历史首创精神，从群众中汲取无穷的智慧和力量。回答第二问必须有比较强归纳能力，即要能把材料2提供党的十八大以来理论和实践创新突破，与成功推进和拓展了中国式现代化结合起来。就是要归纳出：第一，新时代十年来，党对中国式现代化，在认识上不断深化，创立了习近平新时代中国特色社会主义思想，实现了马克思主义中国化时代化新的飞跃，为中国式现代化提供了根本遵循；对中国式现代化深远意义、科学内涵、本质要求及目标蓝图、发展规划、战略导向等方面都作了丰富和完善，初步构建中国式现代化理论体系，使中国式现代化更加清晰、更加科学、更加可感可行，成功推进和拓展了中国式现代化。第二，在实践上不断丰富，推进一系列变革性实践、实现一系列突破性进展、取得一系列标志性成果，为中国式现代化提供了更为完善制度保证、更为坚实的物质基础、更为主动的精神力量，为全面建成社会主义现代化国家奠定了坚实基础和创造了有利条件。当然，现在，这个问题的答案更加明确了，习近平在学习贯彻党二十大精神研讨班开班式上的讲话指出，党的十八大以来，我们党在已有基础上继续前进，不断实现理论和实践上创新突破，成功推进和拓展了中国式现代化。第一，我们在认识上不断深化，创立了习近平新时代中国特色社会主义思想，实现了马克思主义中国化时代化新的飞跃，为中国式现代化提供了根本遵循。我们进一步深化对中国式现代化的内涵和本质认识，概括形成中国式现代化中国特色、本质要求和重大原则，初步构建中国式现代化的理论体系，使中国式现代化更加清晰、更加科学、更加可感可行。第二，我们在战略上不断完善，深入实施科教兴国战略、人才强国战略、乡村振兴战略等一系列重大战略，为中国式现代化提供坚实战略支撑。第三，我们在实践上不断丰富，推进一系列变革性实践、实现一系列突破性进展、取得一系列标志性成果，推动党和国家事业取得历史性成就、发生历史性变革，特别是消除了绝对贫困问题，全面建成小康社会，为中国式现代化提供了更为完善的制度保证、更为坚实的物质基础、更为主动的精神力量。


\subsection*{【名师点拨】}

这里顺便为考生提供一个答题的技巧。由于题目中已经点明了 “结合党的十八大以来的理论和实践创新突破”，因此，考生在作答时，至少应该从理论角度答出习近平新时代中国特色社会主义思想（这是党的十八大以来理论上的最大创新突破），党的二十大对中国式现代化的内涵和本质的认识，以及“三新一高”（新发展阶段、新发展理念、新发展格局、高质量发展），加上题干材料最后已经提示两步走战略安排；从实践角度答出党的十八大以来推进一系列变革性实践、实现一系列突破性进展、取得一系列标志性成果，为中国式现代化提供了更为完善的制度保证、更为坚实的物质基础、更为主动的精神力量。只要回答了这些，就可以拿到大部分的分数。


\subsection*{36. 【答案】}

（1）当时党面临的任务是彻底打败日本侵略者，建立新中国。党只有克服了当时党内存在的主观主义、教条主义、宗派主义的错误倾向，团结一致，才能胜利完成历史任务。只有看齐，才能使全党认识一致、目标明确、力量强大。毛泽东在党的七大预备会议上强调看齐意识，确立了看齐的原则。党的七大使全党在思想上、政治上、组织上达到空前统一和团结，为党后来不断从胜利走向胜利指明了正确方向、开辟了正确道路。


（2）中华民族是有着伟大团结精神的民族，中国共产党是高度集中统一马克思主义政党，思想上的统一、政治上的团结、行动上的一致，是党的事业不断发展壮大的根本所在。实现中华民族伟大复兴必须持续团结奋斗。100多年来，党领导人民团结奋斗，创造了新民主主义革命約伟大成就、社会主义革命和建设的伟大成就、改革开放和社会主义现代化建设的伟大成就、新时代中国特色社会主义的伟大成就。党的历史经验证明，只要全党步调一致、团结统一，就能战胜一切艰难险阻和强大敌人；反之，党和国家事业就会遭受挫折。


\subsection*{【考点点拨】}

结合党的七大召开时党面临的形势和任务，说明毛泽东强调“看齐意识”的重要性；100年来，党和人民取得的一切成就都是团结奋斗的结果


\subsection*{【试题简析】}

本题的第一问有一定的难度，关键是对提问中的“当时”，要有准确的理解，否则就无法回答“当时的形势和党的任务”，当然就更不可能进一步回答毛泽东强调“有了偏差，就喊看齐”的重要性。有一些考生就是因为不了解党的七大预备会议与党的七大的关系，也就把握不好第一问提问所说的“结合当时的形势和党的任务”其实，党的七大预备会议是党的七大的组成部分，因此，“结合当时的形势和党的任务”也就是结合党的七大召开时党面临的形势和任务。明确了这一点，第一问就不难回答了。第一，形势抗日战争即将取得胜利；中国共产党领导的军队和革命力量有了很大的发展，尤其是通过延安整风和党的六届七中全会原则通过的《关于若干历史问题的决议》，全党特别是党的高级干部对中国革命基本问题的认识达到了一致，增强了全党团结。第二，任务——“放手发动群众，壮大人民力量，在我党的领导下，打败日本侵略者，解放全国人民，建立一个新民主主义的中国”。第三，看齐的重要性。党只有克服当时党内存在的主观主义、教条主义、宗派主义的错误倾向，团结一致，才能胜利完成历史任务。要团结一致，就必须向中央基准看齐。只有看齐，才能使全党认识一致、目标明确、力量强大。因此，毛泽东在党的七大预备会议上强调看齐意识，确立了看齐的原则。党的七大以“团结的大会，胜利的大会”载人党的史册，使全党在思想上、政治上、组织上达到空前统一和团结，为党后来不断从胜利走向胜利指明了正确方向、开辟了正确道路。第二问是比较容易回答的，它是党的二十大报告提出的“五个必由之路”的内容之一（“团结奋斗是中国人民创造历史伟业的必由之路”），也是我在预测卷中一再强调的。只要结合已学知识，把材料2、材料3中的基本观点加以归纳，从两个角度展开回答就可以了。首先，要指出中华民族和中国共产党是有着伟大团结精神的。中华民族是有着伟大团结精神的民族，中国共产党是高度集中统一马克思主义政党，思想上的统一、政治上的团结、行动上的一致，是党的事业不断发展壮大的根本所在。实现中华民族伟大复兴必须持续团结奋斗。其次，结合党的历史经验，指出100多年来，党领导人民团结奋斗，创造了新民主主义革命的伟大成就、社会主义革命和建设的伟大成就、改革开放和社会主义现代化建设的伟大成就、新时代中国特色社会主义的伟大成就。党的历史经验证明，只要全党步调一致、团结统一，就能战胜一切艰难险阻和强大敌人；反之，党和国家事业就会遭受挫折。


\subsection*{37. 【答案】}

（1）爱国主义是具体的、历史的，体现了人们对自己祖国的深厚情感，揭示了个人对祖国的依存关系，是人们对自己家园以及民族和文化的归属感、认同感、尊严感与荣誉感的统一，是一个人立德之源、立功之本。爱国不需要理由，是每个人都应当自觉履行的责任和义务。科学家精神是科技工作者在长期科学实践中积累的宝贵精神财富，爱国主义是科学家精神的底色。科学研究虽然没有国界，但是科学家是有祖国的，科学家生存、发展一切物质条件获取和精神慰藉首先得益于祖国。国家是科学家成长发展的坚实依托，失去国家庇佑和保护，科学家也将失去成长和发展的最坚实依托。新时代中国特色社会主义事业是一项前无古人的创造性事业，以爱国主义为底色的科学家精神作为兴国强国之魂的价值和意义将更为凸显。


（2）创新能力是当今国际竞争新优势的集中体现。钱七虎持之以恒为国铸盾的创新实践，生动诠释了科学家勇攀高峰、敢为人先、创新创造的精神内涵。青年时期是创新创造的宝贵时期，当代青年要发扬勇攀高峰、敢为人先的创新精神，不断增强创新意识，夯实创新本领，要在解决国家重大原创性科学问题、突破制约发展的关键核心技术以及勇闯创新“无人区”中不断积累经验、取得成果、演绎精彩。


\subsection*{【考点点拨】}

爱国主义，科学家精神，当代青年如何进行创新创造


\subsection*{【试题简析】}

本题第一问考查爱国主义与科学家精神。科学没有国界，但是科学家有祖国。科学家精神是科技工作者在长期科学实践中积累的宝贵精神财富，中国科学家精神中，摆在首位的就是胸怀祖国、服务人民的爱国精神。爱国主义是中国科学家的光荣传统，也是科学家精神的底色。科学家的生存、发展的一切物质条件获取和精神慰藉首先得益于祖国。国家是科学家成长发展的坚实依托，失去国家庇佑和保护，科学家也将失去成长和发展的最坚实依托。在中华民族伟大复兴的征程上，一代又一代科学家心系祖国和人民，不畏艰难，无私奉献，为科学技术进步、人民生活改善、中华民族发展作出了重大贡献。爱国主义是具体的、历史的，体现了人们对自己祖国的深厚情感，揭示了个人对祖国的依存关系，是人们对自己家园以及民族和文化的归属感、认同感、尊严感与荣誉感的统一，是一个人立德之源、立功之本。爱国不需要理由，是每个人都应当自觉履行的责任和义务。爱国主义是中华民族的民族心、民族魂，是中华民族最重要的精神财富，深深植根于中国人民心中，维系着中华大地上各个民族的团结统一，激励着一代又一代中华儿女为祖国发展繁荣而自强不息、不懈奋斗。中国特色社会主义进入新时代，实现中华民族伟大复兴的中国梦是新时代爱国主义的鲜明主题。新时代中国特色社会主义事业是一项前无古人的创造性事业，以爱国主义为底色的科学家精神作为兴国强国之魂的价值和意义将更为凸显。实现中华民族伟大复兴，更需要继续发扬以爱国主义为底色的科学家精神。第二问考查的实际上是青年人如何进行创新创造的问题。这一问属于开放性问题，即便考生没有提前准备，也可以结合材料（主要是材料2最后一段）进行扩展。比如，将“创新决定未来，创造引领未来”换个说法，就是“创新能力是当今国际竞争新优势的集中体现”。再比如，既然提问是“钱七虎为国铸盾的科学实践”对当代青年有何启示，那么就可以把钱七虎的创新实践再用自己的语言描述一遍。针对此类开放性问题，考生一定要结合题目多写多答，不要怕重复、不要怕“抄材料”，答案要达到一定的字数才有可能拿高分。


\subsection*{【名师点拨】}

爱国主义作为分析题考点，多年来没有考查过。所以，我的2023 版预测卷把它作为重点问题。果然不出所料，2023年的第37题第一问考查的就是爱国主义。


\subsection*{38. 【答案】}

（1）当前，世界之变、时代之变、历史之变正以前所未有的方式展开。一方面，和平、发展、合作、共赢的历史潮流不可阻挡，人心所向、大势所趋决定了人类前途终归光明。另一方面，国际社会面临着日益凸显的传统和非传统安全挑战；恃强凌弱、巧取豪夺、零和博弈等霸权霸道霸凌行径危害深重，和平赤字、发展赤字、安全赤字、治理赤字加重，人类社会面临前所未有的挑战；和平与发展的时代主题能否顶住逆风逆流冲击继续延伸展开，关系全人类的前途命运。


（2） 中国外交政策的宗旨，是维护世界和平、促进共同发展。中国不仅在资金、物资、技术等方面给予发展中国家无私帮助，而且努力维护广大发展中国家的正当权益，提高发展中国家的国际地位，充分体现出中国推动构建人类命运共同体、推动国际关系民主化、弘扬全人类共同价值、正确义利观、全球安全观、全球治理观等理念主张。中国提出共建“一带一路”、全球发展、全球安全等重大倡议，推动创设亚洲基础设施投资银行等多边金融机构，为国际公共产品和国际合作搭建起广阔平台。中国以实际行动为解决当代人类面临的难题贡献着中国智慧；为促进世界和平发展提供了中国方案、注入了新的动力；为发展中国家指明了切实可行的路径，体现了中国的责任与担当，让世界看到了大国应有的样子。


\subsection*{【考点点拨】}

当今世界局势，中国的贡献及外交理念与大国担当


\subsection*{【试题简析】}

第一问有一定难度，关键是要准确审题，明确题干问的“为什么说‘世界又一次站在历史的十字路口’”，实际上是要回答如何认识当今世界的局势的问题。另外，还要明白，所谓“十字路口”，就是说当今世界局势具有两重性（一方面、另一方面），只要明白了这些，作答并不难。当前，世界之变、时代之变、历史之变正以前所未有的方式展开。一方面，和平、发展、合作、共赢的历史潮流不可阻挡，人心所向、大势所趋决定了人类前途终归光明。另一方面，国际社会面临着日益凸显的传统和非传统安全挑战；恃强凌弱、巧取豪夺、零和博弈等霸权霸道霸凌行径危害深重，和平赤字、发展赤字、安全赤字、治理赤字加重，人类社会面临前所未有的挑战；和平与发展的时代主题能否顶住逆风逆流冲击继续延伸展开，关系全人类的前途命运。第二问并不难，提问的前半句是中国作出的巨大贡献，可以结合材料2第一段来扩展：对于发展中国家，中国不仅在资金、物资、技术等方面给予无私帮助，而且努力维护广大发展中国家的正当权益，提高发展中国家的国际地位。提问的后半句是体现的中国外交理念与大国担当，要把具体的理念展开写出来，即推动构建人类命运共同体、推动国际关系民主化、弘扬全人类共同价值、正确义利观、全球安全观、全球治理观等理念主张，再加上材料中已经提及的“一带一路”倡议、全球发展倡议，就形成了整个答案的主体。


\subsection*{【名师点拨】}

我在预测卷中已经明确给考生指出了中国外交的政策宗旨和大国担当，从背诵准备上来说是充足的，关键在于语言组织。从答题技巧上来讲，将提问进行二次改写是一种稳妥且简便易行的方法。也就是说，提问是“‘A’如何体现了‘B’”，那么自然可以回答“什么样的‘A’体现了怎样的“B’”。

\end{document}
